\makeatletter
\def\input@path{{./}{../}{../../}{../../../}{../../../../}{preamble/}{../preamble/}{../../preamble/}{../../../preamble/}{../../../../book/preamble/}}
\makeatother
\input{preamble_exams.tex}
\input{graphs/highlights.tex}
% \GTenHighlightsfalse % Uncomment to disable instructional highlights.
% Population/model, sample evidence, design threshold, and result roles.
\newcommand{\Pop}[1]{\HighlightStroke{GTenHLParameter}{#1}}
\newcommand{\Obs}[1]{\HighlightStroke{GTenHLValue}{#1}}
\newcommand{\Design}[1]{\HighlightStroke{GTenHLThird}{#1}}
\newcommand{\Result}[1]{\HighlightStroke{GTenHLOutcome}{#1}}

\title{Grade 11 Global Economics\\Term 3 --- Lesson 1\\Home Practice C9: Automating and validating hypothesis tests}
\author{}\date{}
\hypersetup{pdftitle={Grade 11 Global Economics - Term 3 - Lesson 1 - C9},pdfsubject={Automating and validating hypothesis tests}}
\begin{document}
\maketitle
\setcounter{tocdepth}{1}
\setcounter{secnumdepth}{2}
\phantomsection\label{toc}
\tableofcontents
\section*{Evaluated criteria}
\begin{description}
\item[\textbf{C9}] Applies computational thinking and development skills to automate hypothesis testing procedures, validate statistical calculations, represent results visually, and generate clear contextual interpretations.
\end{description}
\textbf{How to practise.} Trace the code on paper first. When working independently, run it in a Python environment with NumPy, SciPy, and Matplotlib, then complete the validation tasks before trusting a report. Attempt each task before opening its solution link. Use the return arrows to revisit the problem, group, or contents. All datasets and scenarios in this activity are illustrative.
\clearpage
\section{Introduction}\label{sec:introduction}
\subsection*{Highlight legend}
\phantomsection\label{tab:highlight-legend}
\begin{tabularx}{\linewidth}{@{}>{\raggedright\arraybackslash}p{.22\linewidth}>{\raggedright\arraybackslash}p{.25\linewidth}>{\raggedright\arraybackslash}X@{}}
\toprule
\textbf{Symbol or notation}&\textbf{Concept}&\textbf{Brief meaning}\\\midrule
$\Pop{\theta},\Pop{\theta_0}$ & \Pop{Parameter model} & Population target and null. \\
$\Obs{T},\Obs{SE}$ & \Obs{Calculated evidence} & Statistic and uncertainty. \\
$\Design{\alpha},\Design{c}$ & \Design{Decision configuration} & Level and critical value. \\
$\Result{\text{p-value}}$ & \Result{Tail probability} & Extremeness under the null. \\
\texttt{kind}, \texttt{alt} & \Design{Input choices} & Design and claim direction. \\
\bottomrule
\end{tabularx}
\subsection{Introductory example 1: A mean-testing workflow}\label{ex:intro-one}
A fee analyst has independent observations in USD, a benchmark $\mu_0=20$, and no known population variance. Select a one-sample t model rather than Z solely because data are stored in software. Input parsing, model selection, calculations, and reporting are separate steps.
\subsection{Introductory example 2: A proportion-testing workflow}\label{ex:intro-two}
A lender counts one-year defaults in a random independent sample and asks whether population $p$ exceeds $p_0=0.10$. Before a proportion Z-test at $\alpha=0.05$, software should require integer counts, $0\le x\le n$, and null expected counts at least 10. It still needs a human assessment of sampling and dependence.
\section{Notation Reference and Practice Procedure}\label{sec:notation-and-procedure}
\subsection{Notation reference}\label{tab:notation-reference}
\begin{tabularx}{\linewidth}{@{}>{\raggedright\arraybackslash}p{.16\linewidth}>{\raggedright\arraybackslash}p{.18\linewidth}>{\raggedright\arraybackslash}p{.28\linewidth}>{\raggedright\arraybackslash}X@{}}
\toprule
\textbf{Symbol or notation}&\textbf{Concept}&\textbf{Definition}&\textbf{Meaning in the introductory examples}\\\midrule
$\Pop{\theta},\Pop{\theta_0}$ & \Pop{Parameter model} & $\theta$ can be a mean, proportion, mean paired difference, or independent mean contrast. & \textbullet\ \hyperref[ex:intro-one]{Example 1}: population mean fee.\par\textbullet\ \hyperref[ex:intro-two]{Example 2}: population default rate. \\
$\Obs{T},\Obs{SE}$ & \Obs{Calculated evidence} & $T=(\widehat\theta-\theta_0)/SE$; the selected model determines the standard error and null reference. & \textbullet\ \hyperref[ex:intro-one]{Example 1}: use estimated sample spread with a t reference. \\
$\Design{\alpha},\Design{c}$ & \Design{Decision configuration} & $\alpha$ is chosen in advance; $c$ is a null-distribution quantile for that level and tail. & \textbullet\ \hyperref[ex:intro-two]{Example 2}: level 0.05, right-tail boundary 1.645. \\
$\Result{\text{p-value}}$ & \Result{Tail probability} & Probability of a statistic at least as extreme under the null model, using the planned alternative. & \textbullet\ \hyperref[ex:intro-two]{Example 2}: shade the right tail beyond observed Z. \\
\texttt{kind}, \texttt{alt} & \Design{Input choices} & \texttt{kind} routes the model; \texttt{alt} is \texttt{less}, \texttt{greater}, or \texttt{two-sided}. & \textbullet\ \hyperref[ex:intro-one]{Example 1}: \texttt{mean}.\par\textbullet\ \hyperref[ex:intro-two]{Example 2}: \texttt{proportion}. \\
\bottomrule
\end{tabularx}
\subsection{Practice procedure}
\begin{enumerate}
\item Translate the question into parameter, benchmark, units, study design, alternative, and planned level. Record whether a population SD is genuinely known.
\item Assess the sampling and distribution assumptions outside the program: software cannot establish random sampling, meaningful matching, or independence from a list of numbers.
\item Parse inputs, reject invalid values, select the model, and calculate without rounding intermediate quantities.
\item Validate calculations against a hand-worked reference and a separately implemented trusted routine. Use invariant and invalid-input cases, not only one successful run.
\item Plot the null reference, observed statistic, critical boundary, and p-value region with labels. Distinguish the p-value region from the rejection region.
\item Report hypotheses, assumptions, estimate, standard error, statistic, degrees of freedom, critical values, p-value, level, decision, contextual evidence, and economic importance.
\end{enumerate}
\section{General Worked Problems}\label{sec:general-worked}
\SectionProblemsTableOfContents{\SectionProblemLink{prob:automation}{A complete reusable Python workflow} \\
\SectionProblemLink{prob:validation}{Independent practice: validation cases}}
\subsection{A complete reusable Python workflow}\label{prob:automation}
An analyst samples five independent processing fees, 18, 19, 20, 21, and 22 USD, from an approximately normal population without serious outliers. Population variance is unknown. The preplanned claim is a population mean above 18 USD at $\alpha=0.05$. Supplied reference for $df=4$: right critical 2.132 and $P(T_4\ge2.828)\approx0.0237$. The following program also supports proportions, paired data, and independent groups. It uses the dependencies NumPy, SciPy, and Matplotlib.
\QuestionsTasks[sec:general-worked]{{9}/{sol:automation}/{Explain the decomposition, trace the mean example, identify model assumptions, and describe the report and plot that the program is intended to produce.}}
\subsubsection*{Worked solution}\label{sol:automation}
The parameter is $\Pop{\mu}$, population mean processing fee. Write $H_0:\mu=18$, $H_a:\mu>18$. The independent random sample and approximate normal shape support a small-sample t-test. Unknown variance requires estimating spread with $s$, not entering that estimate as known \texttt{sigma}. The intended hand calculation is
\[
\bar x=20,\quad s^2=\frac{4+1+0+1+4}{4}=2.5,\quad
SE=\frac{\sqrt{2.5}}{\sqrt5}=\sqrt{0.5},\quad
t=\frac{20-18}{\sqrt{0.5}}\approx2.828,\quad df=4.
\]
The supplied reference gives p-value about 0.0237 and critical value 2.132. Therefore \Result{reject $H_0$} at 0.05, supporting a population mean above 18 USD. The estimated difference is 2 USD; a business decision still requires a cost assessment.

The program separates parsing, formula calculation, common tail logic, design routing, reporting, and plotting. \texttt{ddof=1} estimates sample variance with denominator $n-1$. It compares unrounded p-values to the level. These are intended teaching results, not an executed transcript.

\begin{lstlisting}[language=Python,basicstyle=\ttfamily\small,breaklines=true,columns=fullflexible,keepspaces=true]
import math
import numpy as np
from scipy import stats
import matplotlib.pyplot as plt

def numbers(values):
    # Accept a comma-separated string or a sequence.
    if isinstance(values, str):
        values = values.split(",")
    x = np.asarray([float(v) for v in values], dtype=float)
    if x.ndim != 1 or len(x) < 2:
        raise ValueError("Supply at least two observations.")
    if not np.isfinite(x).all():
        raise ValueError("Observations must be finite.")
    return x

def finite(value, name):
    value = float(value)
    if not math.isfinite(value):
        raise ValueError(name + " must be finite.")
    return value

def finish(estimate, null, se, ref, model, df, alt, alpha):
    alpha = finite(alpha, "alpha")
    null = finite(null, "null")
    if not 0 < alpha < 1:
        raise ValueError("alpha must lie strictly between 0 and 1.")
    if not math.isfinite(se) or se <= 0:
        raise ValueError("Standard error must be positive.")
    stat = (estimate - null) / se
    if alt == "greater":
        p = ref.sf(stat)
        critical = (ref.ppf(1 - alpha),)
    elif alt == "less":
        p = ref.cdf(stat)
        critical = (ref.ppf(alpha),)
    elif alt == "two-sided":
        p = min(1.0, 2 * ref.sf(abs(stat)))
        critical = (ref.ppf(alpha / 2),
                    ref.ppf(1 - alpha / 2))
    else:
        raise ValueError("Use less, greater, or two-sided.")
    return dict(estimate=float(estimate), null=null, se=se,
                statistic=float(stat), pvalue=float(p),
                alpha=alpha, critical=critical, df=df,
                model=model, reference=ref, alternative=alt,
                reject=bool(p <= alpha))

def mean_test(data, null, alt, alpha, sigma=None):
    x = numbers(data)
    n = len(x)
    if sigma is not None:
        sigma = finite(sigma, "population sigma")
        if sigma <= 0:
            raise ValueError("Known population sigma must be positive.")
        return finish(x.mean(), null, sigma / math.sqrt(n),
                      stats.norm, "mean Z", None, alt, alpha)
    se = x.std(ddof=1) / math.sqrt(n)
    return finish(x.mean(), null, se, stats.t(df=n-1),
                  "one-sample t", n-1, alt, alpha)

def proportion_test(count, n, p0, alt, alpha):
    count = finite(count, "count")
    n = finite(n, "n")
    p0 = finite(p0, "p0")
    if n <= 0 or not n.is_integer():
        raise ValueError("n must be a positive integer.")
    if not count.is_integer() or not 0 <= count <= n:
        raise ValueError("count must be an integer from 0 to n.")
    if not 0 < p0 < 1:
        raise ValueError("p0 must be strictly between 0 and 1.")
    if min(n*p0, n*(1-p0)) < 10:
        raise ValueError("Null expected counts too small for Z.")
    se = math.sqrt(p0 * (1-p0) / n)
    return finish(count/n, p0, se, stats.norm,
                  "proportion Z", None, alt, alpha)

def paired_test(before, after, null, alt, alpha):
    before, after = numbers(before), numbers(after)
    if len(before) != len(after):
        raise ValueError("Paired lists must have equal length.")
    # Pair by verified unit ID before calling this function.
    result = mean_test(before-after, null, alt, alpha)
    result["model"] = "paired t (before minus after)"
    return result

def welch_test(group1, group2, null, alt, alpha):
    x, y = numbers(group1), numbers(group2)
    n1, n2 = len(x), len(y)
    a = x.var(ddof=1) / n1
    b = y.var(ddof=1) / n2
    if a + b <= 0:
        raise ValueError("Cannot use zero combined variability.")
    df = (a+b)**2 / (a*a/(n1-1) + b*b/(n2-1))
    return finish(x.mean()-y.mean(), null, math.sqrt(a+b),
                  stats.t(df=df), "Welch t", df, alt, alpha)

def run_test(kind, **inputs):
    routes = {"mean": mean_test, "proportion": proportion_test,
              "paired": paired_test, "independent": welch_test}
    if kind not in routes:
        raise ValueError("Unknown study design.")
    # Missing or unexpected argument names raise TypeError.
    return routes[kind](**inputs)

def report(result, parameter, units, claim):
    decision = ("reject H0" if result["reject"]
                else "fail to reject H0")
    evidence = ("Evidence supports " if result["reject"]
                else "Insufficient evidence for ")
    return (
        f"Model: {result['model']}; parameter: {parameter}.\n"
        f"Estimate={result['estimate']:.3f} {units}; "
        f"null={result['null']:.3f} {units}.\n"
        f"SE={result['se']:.3f} {units}; "
        f"statistic={result['statistic']:.3f}; "
        f"df={result['df']}.\n"
        f"Critical boundary/boundaries={result['critical']}; "
        f"p-value={result['pvalue']:.4f}; "
        f"alpha={result['alpha']:.3f}: {decision}.\n"
        + evidence + claim + ".\n"
        "Assess economic size and study assumptions separately."
    )

def plot_null(result):
    ref = result["reference"]
    observed = result["statistic"]
    lo = min(ref.ppf(0.001), observed-1)
    hi = max(ref.ppf(0.999), observed+1)
    grid = np.linspace(lo, hi, 600)
    density = ref.pdf(grid)
    alt = result["alternative"]
    if alt == "greater":
        extreme = grid >= observed
    elif alt == "less":
        extreme = grid <= observed
    else:
        extreme = np.abs(grid) >= abs(observed)
    fig, ax = plt.subplots()
    ax.plot(grid, density, label="Null reference density")
    ax.fill_between(grid, 0, density, where=extreme,
                    alpha=0.3, label="p-value region")
    ax.axvline(observed, linestyle="-.", label="Observed statistic")
    for i, c in enumerate(result["critical"]):
        ax.axvline(c, linestyle="--",
                   label="Critical boundary" if i == 0 else None)
    ax.set(xlabel="Standardized test statistic",
           ylabel="Density", title=result["model"])
    ax.legend()
    plt.show()

# Complete illustrative example. Plan before inspecting outcomes:
# Independent random fees from an approximately normal population,
# no serious outliers; population variance unknown.
# Research claim: population mean fee exceeds 18 USD.
try:
    raw_fees = "18,19,20,21,22"
    result = run_test("mean", data=raw_fees, null=18,
                      alt="greater", alpha=0.05)
    print(report(result, "population mean processing fee",
                 "USD", "a population mean fee above 18 USD"))
    plot_null(result)
except (ValueError, TypeError) as problem:
    print("Input or model error:", problem)
\end{lstlisting}

\textbf{Selection logic.} Use \texttt{mean} with \texttt{sigma} only for genuinely known population SD; omit it for a one-sample t-test. Use \texttt{proportion} for a binary count, \texttt{paired} for correctly aligned before and after lists, and \texttt{independent} for unrelated groups using Welch. Supply all required named inputs; do not silently change a design to avoid an error.

\textbf{Other formulas implemented.} For proportions, $SE_0=\sqrt{p_0(1-p_0)/n}$. For paired data, form $d_i=\text{before}_i-\text{after}_i$ and perform a one-sample t-test on differences with $df=n-1$. For Welch, put $a=s_1^2/n_1$, $b=s_2^2/n_2$; then $SE=\sqrt{a+b}$ and
\[
\nu=\frac{(a+b)^2}{a^2/(n_1-1)+b^2/(n_2-1)}.
\]
Independent groups need within- and between-group independence and suitable shapes. Paired analysis needs independent pairs and suitable differences. The program checks numerical validity, not these substantive conditions.

\textbf{Visual reading.} The curve is the null t density in the worked example. The solid-pattern marker identifies observed $t$; dashed lines show critical boundaries. Shading beyond the observed statistic depicts the p-value, whereas the rejection region begins at the critical boundary. Density height is not itself a probability; area is. The finite plotting window and grid approximate the picture; distribution functions calculate the p-value. Labels and line patterns retain meaning without color.
\secInicio[prob:automation][sec:general-worked]
\subsection{Independent practice: validation cases}\label{prob:validation}
Use the program only after tracing its inputs. For the fee example, the hand reference above is available. Prepare cases that (a) add 100 USD to every fee and to the benchmark; (b) multiply fees and benchmark by 10; (c) reverse a paired difference and its alternative; (d) supply a negative known SD, a non-integer default count, a missing observation, or a zero-variation unknown-variance sample. Also compare the mean example with \texttt{scipy.stats.ttest\_1samp} using the same alternative.
\QuestionsTasks[sec:general-worked]{{9}/{sol:validation}/{Predict invariant results, specify expected errors and tolerances, then explain how students should compare independent calculations. Do not infer validity solely from a matching decision.}}
\subsubsection*{Worked solution}\label{sol:validation}
For (a), the mean and null both increase by 100 while their difference and $SE$ are unchanged, so statistic, p-value, and decision are unchanged. For (b), numerator and standard error both multiply by 10, so those same testing quantities remain unchanged. Units and reported estimates must change appropriately. For (c), consistent reversal changes the statistic's sign and swaps left with right; p-value and decision are unchanged.

Invalid SD, non-integer count, non-finite values, and zero estimated standard error should produce clear errors. Blank entries fail numerical conversion rather than being silently removed; missing matched observations require an explicit missing-data plan. Verify pair alignment by IDs before numerical calls.

For the mean reference, independently calculate $\bar x$, $s$, $SE$, $t$, and $df=4$. Compare unrounded formula output with a separate \texttt{ttest\_1samp} call using \texttt{alternative="greater"}. A tolerance such as $10^{-10}$ is suitable for comparing two software calculations on these small finite inputs; compare the rounded hand values only to their shown precision. Compare statistic and p-value separately, and inspect hypotheses, assumptions, and tail configuration. Agreement on a reject decision alone could hide the same formula mistake in both programs. Students should record their own executions and discrepancies; no program output is supplied as an executed result here.
\secInicio[prob:validation][sec:general-worked]
\section{Advanced Problems}\label{sec:advanced}
\SectionProblemsTableOfContents{\SectionProblemLink{prob:debug}{Debug, extend, and explain}}
\subsection{Debug, extend, and explain}\label{prob:debug}
A draft routine computes a proportion standard error with $\hat p$ instead of $p_0$, uses \texttt{1-cdf(t)} for every alternative, calculates sample SD with denominator $n$, and reports that a p-value of 0.02 means the null is 2\% likely. Its plot shades beyond the critical value and labels that area as the p-value. The manager asks for batch processing of many economic reports.
\QuestionsTasks[sec:advanced]{{9}/{sol:debug}/{Explain every defect and design an extension with input handling, validation, visual output, and contextual reporting.}}
\subsubsection*{Worked solution}\label{sol:debug}
The null proportion test needs $p_0(1-p_0)/n$, because the reference distribution is computed under $H_0$. The p-value depends on the preplanned alternative: use the lower cumulative probability for left tails, the survival probability for right tails, and twice the survival probability at the absolute statistic for a symmetric two-sided reference. A survival function also avoids subtracting a probability very close to one. Sample variance uses $n-1$, implemented by \texttt{ddof=1}. A p-value is conditional on the null model, not the probability of its truth. Shading beyond a critical boundary depicts the rejection region; p-value shading begins at the observed extremeness boundary.

A batch extension should read a table with explicit design, benchmark, units, alternative, level, dataset identifier, and assumption notes. Validate each row; retain errors alongside the row ID rather than substituting a default test. For paired rows require verified unit IDs; for independent rows confirm no overlap. Route each row through the reusable functions and store estimate, $SE$, statistic, $df$, critical boundaries, p-value, decision, and assumption limitations. Add an effect plot in original economic units alongside each labeled null-distribution plot. Include the hypotheses and research claim in the exported report, and describe practical significance rather than reporting only a Boolean result.

Validate the extension with a hand reference, a trusted separate implementation, translation/scaling invariants, swapped groups with a consistently reversed alternative, and malformed rows. For two-sided Welch tests, swapping groups reverses the statistic while retaining its absolute magnitude, degrees of freedom, and p-value. If many claims are tested, decide how to address multiple testing before interpreting the batch; a collection of unadjusted 5\% tests does not retain a 5\% probability of any false rejection. These are development tasks for students; they require execution and review in their own work, not assumptions of correctness from readable code.
\secInicio[prob:debug][sec:advanced]
\end{document}

